\SourceLeaf{491}%
Skagen angestellt sind, werden verglichen werden müssen, so könnte es dem Zutrauen, welches die Resultate für sich fordern sollen, nachtheilig werden, wenn jene Beobachtungen nicht bloss mit verschiedenen, sondern sogar mit ganz verschiedenartigen Instrumenten gewonnen wären. Am besten wird es daher sein, wenn die hiesigen Beobachtungen nicht bloss am Meridiankreise, sondern überdiess noch am Zenithsector, und zwar an dem \textit{nämlichen} Zenithsector, womit an allen andern Orten observirt wurde, angestellt werden.

Der Zenithsector womit die Dänischen Bestimmungen in Lauenburg, Lysabbel und Skagen, auch in Kopenhagen, in den Jahren 1819 und 1820 gemacht sind, ist eben derselbe, womit bei der englischen Gradmessung im Jahre 1802 vom verst. General \textsc{Mudge} beobachtet ist. Mit eben dem Instrumente wurde im Jahre 1818 gemeinschaftlich von englischen und französischen Gelehrten in Dünkirchen als dem nordlichsten Endpunkt der französischen Gradmessung eine Reihe von Beobachtungen gemacht, und nachher dasselbe vom engl. Gouvernement zu der Dänischen Gradmessung auf 2 Jahre geliehen. Alle diese Umstände machen es um so wünschenswerther, dass dieses Instrument auch bei der Hannoverschen Gradmessung gebraucht werde, und es wird dabei besonders wichtig und lehrreich sein, dasselbe mit dem hiesigen Reichenbachschen Meridiankreise vermittelst gleichzeitiger und in Einem Ort gemachten Beobachtungen zu vergleichen.

Es ist nicht zu zweifeln, dass für die Hannoversche Gradmessung dieses Instrument mit derselben Bereitwilligkeit werde hergeliehen werden, wie für die Dänische. — — — Ich füge noch folgende Umstände bei. Dieser von \textsc{Ramsdon} verfertigte Zenithsector gehört dem Board of Ordnance, und steht in letzter Instanz unter dem Herzog von \textsc{Wellington}, mittelbar aber, nach dem vor kurzem erfolgten Tode des General \textsc{Mudge}, unter dem Oberstlieutenant \textsc{Colby}.

Aus vorstehender Darstellung geht hervor, dass die Zeit des Anfangs der grössern zur Gradmessung gehörigen Operationen, die besondere Vorkehrungen erfordern, in diesem Augenblick sich noch nicht bestimmt angeben lässt. Auf alle Fälle aber wird im Laufe des gegenwärtigen Jahres nichts weiter geschehen können, als ein Theil der astronomischen Beobachtungen auf hiesiger Sternwarte (womit ich auch bereits einen Anfang gemacht habe) und höchstens einige Vorarbeiten in hiesiger Umgegend.

Göttingen den 11. August 1820.
\begin{flushright}C. F. \textsc{Gauss}.\end{flushright}

\par\bigskip
\pdfbookmark[0]{P. M. betreffend die Gegenwart des Hofr. Gauss bei einem Theil der Dänischen Gradmessung.}{pm-gradmessung}
\begin{center}\large P. M. BETREFFEND DIE GEGENWART DES HOFR. GAUSS, BEI EINEM THEIL DER OPERATION DER DÄNISCHEN GRADMESSUNG.\end{center}
\par
Über die in dem königl. Dänischen Gebiete angeordnete und dem Dänischen Astronomen \textsc{Schumacher} übertragene Gradmessung, ihren Zweck, die Vortheile, die man für die Wissenschaften davon zu erwarten berechtigt ist, so wie über die dabei angewandten Hülfsmittel habe ich schon früher die Ehre gehabt, Ew. — — einen unterthänigsten Bericht abzustatten. Es gereicht mir um so mehr zum Vergnügen, dass eben jetzt ein anderer Astronom, der Doctor \textsc{Olbers}. der auch im vorigen Jahre von den Operationen in Lauenburg Augenzeuge gewesen, in der Zeitschrift Allgemeine Geographische Ephemeriden (Bd. 17 Stück 3) mit den meinigen ganz übereinstimmende Ansichten dieser Unternehmung bekannt gemacht hat, da mein eignes Urtheil über die Persönlichkeit des Dänischen Astronomen insofern vielleicht befangen scheinen könnte, weil dieser früher (1809) seine letzte Ausbildung auf hiesiger Universität erhalten hatte.
